\documentclass[11pt]{article}

\usepackage[margin=2.5cm]{geometry}
\usepackage{amsmath,amssymb,amsthm}
\usepackage{bm}

\title{Solutions Sheet -- Calculus}
\author{}
\date{}

\begin{document}
\maketitle
\subsection*{Part A -- Fundamental Exercises}

\paragraph{Exercise 1. Gradients of scalar functions.}

\begin{enumerate}
\item \( f(x,y)=x^2y+3y^2 \).

\[
\nabla f(x,y)=\left(\frac{\partial f}{\partial x},\frac{\partial f}{\partial y}\right)
=(2xy,\;x^2+6y).
\]

\item \( f(x,y,z)=e^{xy}+z^2 \).

\[
\nabla f(x,y,z)=\big(y e^{xy},\; x e^{xy},\; 2z\big).
\]

\item \( f(\bm{x})=\|A\bm{x}\|^2=\bm{x}^T A^T A \bm{x} \).

We recognize a quadratic form. We have
\[
\nabla f(\bm{x})
= 2A^T A\bm{x}.
\]

\end{enumerate}

\bigskip

\paragraph{Exercise 2. Hessian and convexity.}

\( f(x,y)=x^4+2x^2y^2+y^4 = (x^2+y^2)^2. \)

\[
\frac{\partial f}{\partial x}=4x^3+4xy^2=4x(x^2+y^2),\qquad
\frac{\partial f}{\partial y}=4y^3+4x^2y=4y(x^2+y^2).
\]

Second derivatives:
\[
\frac{\partial^2 f}{\partial x^2}=12x^2+4y^2,\quad
\frac{\partial^2 f}{\partial y^2}=4x^2+12y^2,\quad
\frac{\partial^2 f}{\partial x\partial y}
=\frac{\partial^2 f}{\partial y\partial x}=8xy.
\]

Thus
\[
H_f(x,y)=
\begin{pmatrix}
12x^2+4y^2 & 8xy\\[4pt]
8xy & 4x^2+12y^2
\end{pmatrix}.
\]

For convexity: we note that
\[
f(x,y)=(x^2+y^2)^2,
\]
and that \(h(x,y)=x^2+y^2\) is convex, and \(g(r)=r^2\) is convex and increasing on \([0,\infty)\).  
The composition \(f=g\circ h\) is therefore convex.  
One can also check that the Hessian is positive semidefinite (all eigenvalues \(\ge 0\)).

\bigskip

\paragraph{Exercise 3. Jacobian of a vector-valued map.}

\(\Phi(u,v)=(u e^{v},\,u^2+v).\)

\[
J_\Phi(u,v)=
\begin{pmatrix}
\partial_u(ue^v) & \partial_v(ue^v)\\[4pt]
\partial_u(u^2+v) & \partial_v(u^2+v)
\end{pmatrix}
=
\begin{pmatrix}
e^v & u e^v\\[4pt]
2u & 1
\end{pmatrix}.
\]

\[
\det J_\Phi(u,v)= e^v\cdot 1 - 2u\cdot u e^v = e^v(1-2u^2).
\]

\bigskip

\paragraph{Exercise 4. One-dimensional change of variables.}

\[
I=\int_0^1 \frac{1}{\sqrt{1-x^2}}\,dx.
\]

Set \(x=\sin\theta\), with \(\theta\in[0,\pi/2]\).  
Then \(dx=\cos\theta\,d\theta\) and \(\sqrt{1-x^2}=\sqrt{1-\sin^2\theta}=\cos\theta\).

\[
I=\int_{\theta=0}^{\pi/2} \frac{1}{\cos\theta}\cos\theta\,d\theta
=\int_0^{\pi/2} d\theta
=\frac{\pi}{2}.
\]

\bigskip


\paragraph{Exercise 5. Polar coordinates and Jacobian.}

\[
\int_{B(0,1)} (x^2+y^2)\,dx\,dy,
\]
where \(B(0,1)\) is the unit disk.

In polar coordinates \(x=r\cos\theta,y=r\sin\theta\), \(x^2+y^2=r^2\), \(dx\,dy=r\,dr\,d\theta\), and \(0\le r\le 1\), \(0\le\theta\le2\pi\).

\[
\int_0^{2\pi}\int_0^1 r^2 \cdot r\,dr\,d\theta
= \int_0^{2\pi}\int_0^1 r^3\,dr\,d\theta
= \int_0^{2\pi}\left[\frac{r^4}{4}\right]_0^1 d\theta
= \int_0^{2\pi}\frac14\,d\theta
=\frac14\cdot 2\pi=\frac{\pi}{2}.
\]

\bigskip

\paragraph{Exercise 6. Multiple integrals and Fubini.}

\[
\int_0^1\int_0^1 xy\,dx\,dy
= \int_0^1\left[\frac{x^2}{2}y\right]_{x=0}^{1}dy
= \int_0^1 \frac{y}{2}\,dy
= \left[\frac{y^2}{4}\right]_0^1=\frac14.
\]

Conversely,
\[
\int_0^1\int_0^1 xy\,dy\,dx
= \int_0^1\left[\frac{y^2}{2}x\right]_{y=0}^{1}dx
= \int_0^1 \frac{x}{2}\,dx = \frac14.
\]

Both orders give the same value, as guaranteed by Fubini’s theorem.

\bigskip

\paragraph{Exercise 7. Multivariate chain rule.}

\(g(u,v)=u^2+v^2\), \(\Phi(x,y)=(x^2-y,\,xy)\).  
Thus
\[
f(x,y)=g(\Phi(x,y))=(x^2-y)^2+(xy)^2.
\]

\textbf{(a) Direct gradient.}

\[
f_x = 2(x^2-y)\cdot 2x + 2xy\cdot y = 4x(x^2-y)+2xy^2,
\]
\[
f_y = 2(x^2-y)\cdot(-1) + 2xy\cdot x = -2(x^2-y)+2x^2y.
\]

\textbf{(b) Jacobian and chain rule.}

\[
J_\Phi(x,y)=
\begin{pmatrix}
\partial_x(x^2-y) & \partial_y(x^2-y)\\[4pt]
\partial_x(xy) & \partial_y(xy)
\end{pmatrix}
=
\begin{pmatrix}
2x & -1\\[4pt]
y & x
\end{pmatrix}.
\]

\[
\nabla g(u,v)=(2u,2v).
\]

Hence
\[
\nabla f(x,y)=J_\Phi(x,y)^\top\,\nabla g(\Phi(x,y))
=
\begin{pmatrix}
2x & y\\[4pt]
-1 & x
\end{pmatrix}
\begin{pmatrix}
2(x^2-y)\\[4pt]
2xy
\end{pmatrix},
\]
which, after computation, gives back the same components as above (verification left to the students).

\bigskip

\paragraph{Exercise 8. Linear change of variables in \(\mathbb{R}^2\).}

Triangle \(D\) with vertices \((0,0)\), \((1,0)\), \((1,1)\).  
It can be described as
\[
D = \{(x,y): 0\le x\le1,\;0\le y\le x\}.
\]

\textbf{(a) Transformation.}  
We take
\[
x=u,\qquad y=uv,\qquad 0\le u\le1,\;0\le v\le1.
\]
Then \(y\le x\) and we cover exactly \(D\).

\textbf{(b) Jacobian.}

\[
\frac{\partial(x,y)}{\partial(u,v)} =
\begin{pmatrix}
1 & 0\\[4pt]
v & u
\end{pmatrix},
\qquad
\det J = u.
\]

\textbf{(c) Integral via change of variables.}

\[
\int_D (x+y)\,dx\,dy
= \int_0^1\int_0^1 (u+uv)\,u\,dv\,du
= \int_0^1\int_0^1 (u^2+u^2v)\,dv\,du.
\]

\[
\int_0^1 (u^2+u^2v)\,dv
= u^2\left[v+\frac{v^2}{2}\right]_0^1
= u^2\left(1+\frac12\right)
= \frac{3}{2}u^2.
\]

Thus
\[
\int_D (x+y)\,dx\,dy
= \int_0^1 \frac{3}{2}u^2\,du
= \frac{3}{2}\cdot\frac{1}{3}
= \frac{1}{2}.
\]

\textbf{(d) Direct verification.}

\[
\int_0^1\int_0^x (x+y)\,dy\,dx
= \int_0^1\left[xy+\frac{y^2}{2}\right]_{0}^{x}dx
= \int_0^1\left(x^2+\frac{x^2}{2}\right)dx
= \int_0^1 \frac{3}{2}x^2\,dx
= \frac{1}{2}.
\]

\bigskip
\bigskip

\subsection*{Part B -- Advanced and Finance-Oriented Exercises}

\bigskip

\paragraph{Exercise 9. Black--Scholes setup and lognormality.}

We have
\[
S_T = S_0\exp\big((r-\tfrac12\sigma^2)T+\sigma\sqrt{T}Z\big),
\qquad Z\sim N(0,1).
\]

Let
\[
X=\ln S_T = \ln S_0 + (r-\tfrac12\sigma^2)T + \sigma\sqrt{T}Z.
\]
Thus \(X\) is Gaussian:
\[
X\sim \mathcal{N}\big(m,\;s^2\big),
\quad
m = \ln S_0 + (r-\tfrac12\sigma^2)T,\quad s^2=\sigma^2 T.
\]

The density of \(S_T=e^X\) is therefore lognormal:
\[
f_{S_T}(s)=\frac{1}{s\,\sqrt{2\pi\sigma^2T}}
\exp\left(-\frac{(\ln s - m)^2}{2\sigma^2 T}\right),\qquad s>0.
\]

For the expectation:
\[
\mathbb{E}[S_T] = \mathbb{E}[e^X]
= \exp\left(m+\frac{s^2}{2}\right)
= \exp\left(\ln S_0 + (r-\tfrac12\sigma^2)T + \frac{\sigma^2 T}{2}\right)
= S_0 e^{rT}.
\]

\bigskip

\paragraph{Exercise 10. Call and put prices as integrals over \(Z\).}

We write
\[
S_T = S_0\exp\left((r-\tfrac12\sigma^2)T+\sigma\sqrt{T}Z\right),\qquad Z\sim N(0,1).
\]

The event \(\{S_T>K\}\) is equivalent to
\[
Z > z^\ast
:= \frac{\ln(K/S_0) - (r-\tfrac12\sigma^2)T}{\sigma\sqrt{T}}
= -d_2,
\]
where 
\[
d_1 = \frac{\ln(S_0/K) + (r+\tfrac12\sigma^2)T}{\sigma\sqrt{T}},
\qquad
d_2 = d_1 - \sigma\sqrt{T}.
\]

The undiscounted call option:
\[
C = \mathbb{E}[(S_T-K)^+]
= \int_{z^\ast}^{\infty} \big(S_0 e^{(r-\tfrac12\sigma^2)T+\sigma\sqrt{T}z} - K\big)\,\phi(z)\,dz,
\]
where \(\phi\) is the standard normal density.

Similarly, the put:
\[
P = \mathbb{E}[(K-S_T)^+]
= \int_{-\infty}^{z^\ast} \big(K - S_0 e^{(r-\tfrac12\sigma^2)T+\sigma\sqrt{T}z}\big)\,\phi(z)\,dz.
\]

\bigskip

\paragraph{Exercise 11. Vega of a European call option.}

We can start from the usual expression (discounting possibly omitted here):
\[
C(S_0,K,\sigma,T)=S_0\Phi(d_1)-K e^{-rT}\Phi(d_2),
\]
with
\[
d_1 = \frac{\ln(S_0/K) + (r+\tfrac12\sigma^2)T}{\sigma\sqrt{T}},
\qquad
d_2=d_1-\sigma\sqrt{T}.
\]

Differentiation with respect to \(\sigma\) gives
\[
\text{Vega} = \frac{\partial C}{\partial \sigma}
= S_0\phi(d_1)\sqrt{T}.
\]

(By starting from the integral representation in Exercise 14, one can justify exchanging derivative and integral and recover the same formula after simplifications.)

\bigskip

\paragraph{Exercise 12. Vega of a European put and put--call parity.}

The Black--Scholes put:
\[
P(S_0,K,\sigma,T)=K e^{-rT}\Phi(-d_2)-S_0\Phi(-d_1).
\]

Differentiating with respect to \(\sigma\), we obtain
\[
\frac{\partial P}{\partial \sigma}
= K e^{-rT}\phi(d_2)\frac{\partial(-d_2)}{\partial \sigma}
- S_0\phi(d_1)\frac{\partial(-d_1)}{\partial \sigma}.
\]
After computation, we again obtain
\[
\frac{\partial P}{\partial \sigma}
= S_0\phi(d_1)\sqrt{T}.
\]

Alternatively, we use put–call parity:
\[
C-P = S_0 - K e^{-rT}.
\]
Differentiating with respect to \(\sigma\) (the right-hand side does not depend on \(\sigma\)),
\[
\frac{\partial C}{\partial \sigma} - \frac{\partial P}{\partial \sigma} =0
\quad\Rightarrow\quad
\text{Vega}_{\text{call}} = \text{Vega}_{\text{put}}.
\]

\bigskip

\paragraph{Exercise 13. Digital option price as probability.}

Digital call (undiscounted):
\[
D = \mathbb{E}\big[\mathbf{1}_{\{S_T>K\}}\big]
= \mathbb{P}(S_T>K).
\]

In terms of \(Z\):
\[
\mathbb{P}(S_T>K) = \mathbb{P}\left(Z > \frac{\ln(K/S_0)-(r-\tfrac12\sigma^2)T}{\sigma\sqrt{T}}\right)
= \Phi(d_2),
\]
where \(d_2\) is as above (using \(1-\Phi(-d_2)=\Phi(d_2)\)).

We also know (from option theory) that
\[
\frac{\partial C}{\partial K} = -e^{-rT}\mathbb{P}(S_T>K) = -e^{-rT}\Phi(d_2),
\]
which links the price of the digital option to the derivative of the call with respect to \(K\).

\bigskip

\paragraph{Exercise 14. Greeks and Jacobians for a two-instrument portfolio.}

\[
P(\theta)=
\begin{pmatrix}
C(\theta)\\
P(\theta)
\end{pmatrix},
\qquad \theta=(S_0,\sigma).
\]

\textbf{(a) Jacobian.}
\[
J_P(\theta)=
\begin{pmatrix}
\displaystyle \frac{\partial C}{\partial S_0} & \displaystyle \frac{\partial C}{\partial \sigma} \\[8pt]
\displaystyle \frac{\partial P}{\partial S_0} & \displaystyle \frac{\partial P}{\partial \sigma}
\end{pmatrix}.
\]

\textbf{(b) Interpretation.}

\begin{itemize}
\item \(\frac{\partial C}{\partial S_0} = \Delta_{\text{call}}\),
\item \(\frac{\partial C}{\partial \sigma} = \text{Vega}_{\text{call}}\),
\item \(\frac{\partial P}{\partial S_0} = \Delta_{\text{put}}\),
\item \(\frac{\partial P}{\partial \sigma} = \text{Vega}_{\text{put}}\).
\end{itemize}

\textbf{(c) Determinant.}

The determinant
\[
\det J_P(\theta) = 
\Delta_{\text{call}}\cdot \text{Vega}_{\text{put}}
- \Delta_{\text{put}}\cdot \text{Vega}_{\text{call}}
\]
measures (locally) the invertibility of the mapping \(\theta\mapsto (C,P)\).  
A determinant close to zero means that the variations of \((C,P)\) with respect to \((S_0,\sigma)\) are almost colinear: the inversion problem (recovering \((S_0,\sigma)\) from \((C,P)\)) is ill-conditioned, which is unfavorable numerically and for hedging.

\bigskip

\paragraph{Exercise 15. Laplacian, heat equation and Black--Scholes PDE.}

We start from the PDE
\[
\frac{\partial V}{\partial t}
+ \frac{1}{2}\sigma^2 S^2 \frac{\partial^2 V}{\partial S^2}
+ r S \frac{\partial V}{\partial S}
- r V = 0.
\]

We impose in the statement the change of variables
\[
x = \ln\!\left(\frac{S}{K}\right), \qquad
\tau = \frac{1}{2}\sigma^2 (T - t),
\]
and
\[
V(S,t) = K e^{\alpha x + \beta \tau} u(x,\tau),
\]
for suitably chosen constants \(\alpha,\beta\). After computing the derivatives and substituting into the PDE, the terms in \(u_x\) and \(u\) cancel if we choose
\[
\alpha = \frac{1}{2} - \frac{r}{\sigma^2}
\]
and an appropriate \(\beta\) (given in the statement). We then obtain
\[
\frac{\partial u}{\partial \tau} = \frac{\partial^2 u}{\partial x^2},
\]
which is the heat equation.

Since the fundamental solution of the heat equation is Gaussian, the solutions can be expressed via Gaussian integrals, which explains the appearance of normal distributions in the Black–Scholes formula.

\bigskip

\paragraph{Exercise 16. Stability of an ODE using Taylor expansion.}

We consider the ODE
\[
y'(t)=f(y(t)),\quad f\in C^1, \quad f(y^\ast)=0,
\]
and \(z(t)=y(t)-y^\ast\).

\begin{itemize}
\item \(z'(t) = f(y^\ast+z(t))\).

Taylor expansion:
\[
f(y^\ast+z) = f(y^\ast)+f'(y^\ast)z+R(z) = f'(y^\ast)z+R(z),
\]
with \(R(z)=\mathcal{O}(z^2)\).

Approximate (linearized) equation:
\[
z'(t)\approx f'(y^\ast)z(t)=\lambda z(t).
\]

Solution:
\[
z(t)=z(0)e^{\lambda t},\quad \lambda=f'(y^\ast).
\]

\item If \(\lambda<0\), then \(z(t)\to 0\) as \(t\to\infty\): the equilibrium is locally asymptotically stable.  
If \(\lambda>0\), perturbations grow: the equilibrium is unstable.

\item Application to \(y' = y-y^3\).

\(f(y)=y-y^3=y(1-y^2)\).

The equilibria are \(y^\ast\in\{0,1,-1\}\).

\[
f'(y) = 1-3y^2.
\]

Thus
\[
f'(0)=1>0 \Rightarrow y^\ast=0 \text{ unstable},
\]
\[
f'(1)=1-3=-2<0,\quad f'(-1)=-2<0
\Rightarrow y^\ast=\pm1 \text{ locally asymptotically stable}.
\]

\end{itemize}

\end{document}